% Pertemuan 7 Praktikum Pemrograman (IF25-11002). Dihasilkan oleh tools/build.py dari tools/konten/p07.py.
% Kompilasi: xelatex P07.tex (dua kali). Catatan: xelatex -jobname=P07-catatan "\def\catatan{1}% Pertemuan 7 Praktikum Pemrograman (IF25-11002). Dihasilkan oleh tools/build.py dari tools/konten/p07.py.
% Kompilasi: xelatex P07.tex (dua kali). Catatan: xelatex -jobname=P07-catatan "\def\catatan{1}% Pertemuan 7 Praktikum Pemrograman (IF25-11002). Dihasilkan oleh tools/build.py dari tools/konten/p07.py.
% Kompilasi: xelatex P07.tex (dua kali). Catatan: xelatex -jobname=P07-catatan "\def\catatan{1}% Pertemuan 7 Praktikum Pemrograman (IF25-11002). Dihasilkan oleh tools/build.py dari tools/konten/p07.py.
% Kompilasi: xelatex P07.tex (dua kali). Catatan: xelatex -jobname=P07-catatan "\def\catatan{1}\input{P07.tex}"
\documentclass[t]{beamer}
\usetheme{ITERA}
\input{catatan-preamble.tex}
\renewcommand\ITERAmeeting{Pertemuan 7}
\begin{document}
\SlideJudul{IF25-11002 PRAKTIKUM PEMROGRAMAN}{Pertemuan 7\\Review dan Simulasi UTS}{Merangkai materi Pertemuan 1 sampai 6 dan berlatih dengan soal bergaya UTS}{Tim Pengajar}{Tahun 2026. Sejajar dengan Pertemuan 7 Algoritma Pemrograman: Review dan Simulasi UTS.}
\note{Buka dengan menanyakan satu hal dari pertemuan lalu yang masih membingungkan.}

\begin{frame}{Dua mata kuliah, satu minggu}
\framesubtitle{Sebelum mulai}
Topik minggu ini sama di dua mata kuliah, dengan peran yang berbeda.
\vspace{10pt}

\begin{columns}[T]
\begin{column}{0.47\textwidth}
\begin{Blok}{Algoritma: Review dan Simulasi UTS}
Review pseudocode, flowchart, dan trace table untuk materi pertemuan 1 sampai 6.
\end{Blok}
\end{column}
\begin{column}{0.47\textwidth}
\begin{BlokGelap}{Praktikum: Review UTS}
Review C++ pertemuan 1 sampai 6 dan simulasi soal A (kode) dan soal B (tracing dan penjelasan).
\end{BlokGelap}
\end{column}
\end{columns}
\note{Tanyakan apa yang sudah dibahas di Algoritma minggu ini. Gunakan jawabannya sebagai jembatan ke praktikum.}
\end{frame}

\begin{frame}{Saat pulang nanti, kalian bisa}
\framesubtitle{Target hari ini}
\begin{columns}[T]
\begin{column}{0.47\textwidth}
\begin{Langkah}{1}{Menerapkan}
Menyelesaikan masalah yang menggabungkan masukan, ekspresi, percabangan, dan perulangan dalam satu program C++ dalam batas waktu ujian

{\footnotesize\color{ITERAInkMuted}Sub-CPMK 7.1, CPMK0607}
\end{Langkah}
\end{column}
\begin{column}{0.47\textwidth}
\begin{Langkah}{2}{Menjelaskan}
Menelusuri program yang menggabungkan percabangan dan perulangan bersarang, lalu menjelaskan setiap perubahan nilai dengan istilah yang tepat

{\footnotesize\color{ITERAInkMuted}Sub-CPMK 7.2, CPMK0608}
\end{Langkah}
\end{column}
\end{columns}
\vspace{8pt}
\begin{Catatan}
Dua kemampuan ini dinilai sama berat di Exit Ticket, tugas, UTS, dan UAS.
\end{Catatan}
\note{Bacakan target dengan pelan. Kata kuncinya menerapkan dan menjelaskan.}
\end{frame}

\begin{frame}{Ingat minggu lalu}
\framesubtitle{Review, 5 menit}
\begin{columns}[T]
\begin{column}{0.31\textwidth}
\begin{BlokGelap}{Pertanyaan 1}
Apa rumus banyak bintang pada piramida baris b?
\end{BlokGelap}
\end{column}
\begin{column}{0.31\textwidth}
\begin{BlokGelap}{Pertanyaan 2}
Apa beda \texttt{break} dan \texttt{continue} di perulangan bersarang?
\end{BlokGelap}
\end{column}
\begin{column}{0.31\textwidth}
\begin{BlokGelap}{Pertanyaan 3}
Bagaimana menghitung jam dari menit dengan pembulatan ke atas?
\end{BlokGelap}
\end{column}
\end{columns}
\vspace{8pt}
Jawab di kertas dulu, lalu cocokkan dengan teman sebelah.\note{Gunakan tiga pertanyaan ini sebagai diagnosis cepat. Catat topik yang paling banyak salah untuk dibahas di bagian materi.}
\end{frame}

\Bagian{1}{Masalah pembuka}{Satu soal, semua materi}
\note{Masalah pembuka 10 menit.}

\begin{frame}[fragile]{Satu soal, semua materi}
\framesubtitle{Masalah minggu ini}
\begin{columns}[T]
\begin{column}{0.55\textwidth}
{\small Soal UTS jarang menguji satu topik saja. Contohnya tarif parkir: membaca masukan, membulatkan menit ke jam dengan \texttt{/} dan \texttt{\%}, memilih tarif dengan \texttt{if}, dan mengulang untuk beberapa kendaraan.

\vspace{6pt}
Hari ini kita berlatih memecah soal seperti itu menjadi bagian-bagian yang sudah kalian kuasai.\par}
\vspace{6pt}
\begin{Catatan}
\textbf{Diskusi 2 menit.} Materi pertemuan mana saja yang dipakai soal ini? Bagian mana yang menurut kalian paling rawan salah?
\end{Catatan}
\end{column}
\begin{column}{0.41\textwidth}
\begin{Keluaran}[]
M 61   -> 2 jam, tarif 3000
B 180  -> 3 jam, tarif 11000
X 10   -> Jenis tidak dikenal
\end{Keluaran}
\end{column}
\end{columns}
\note{Tampung beberapa jawaban tanpa menilai benar salah.}
\end{frame}

\begin{frame}{Dari langkah ke instruksi}
\framesubtitle{Sambungan dengan Algoritma}
\begin{columns}[T]
\begin{column}{0.31\textwidth}
\begin{Langkah}{1}{Baca soal}
Garis bawahi masukan, keluaran, dan aturan.
\end{Langkah}
\end{column}
\begin{column}{0.31\textwidth}
\begin{Langkah}{2}{Pecah}
Petakan setiap aturan ke konstruksi C++ yang sudah dipelajari.
\end{Langkah}
\end{column}
\begin{column}{0.31\textwidth}
\begin{Langkah}{3}{Uji}
Siapkan data uji, termasuk nilai batas, sebelum menulis kode.
\end{Langkah}
\end{column}
\end{columns}
\vspace{10pt}
Langkah ini sama dengan yang dilatih di Algoritma minggu ini, hanya hasil akhirnya kode C++.\note{Hubungkan eksplisit ke Algoritma minggu ini.}
\end{frame}

\Bagian{2}{Coba dulu, baru dijelaskan}{25 menit eksperimen di GDB Online}
\note{Struggle 25 menit. Berkeliling, jangan menjelaskan materi dulu.}

\begin{frame}{Misi kalian}
\framesubtitle{Struggle, 25 menit}
\begin{columns}[T]
\begin{column}{0.31\textwidth}
\begin{Langkah}{1}{Pilih}
Kerjakan \texttt{handson1-parkir.cpp} tanpa bantuan selama 20 menit.
\end{Langkah}
\end{column}
\begin{column}{0.31\textwidth}
\begin{Langkah}{2}{Uji}
Uji dengan 60 dan 61 menit. Apakah hasilnya berbeda satu jam?
\end{Langkah}
\end{column}
\begin{column}{0.31\textwidth}
\begin{Langkah}{3}{Catat}
Tulis satu hal yang membuat kalian tersendat.
\end{Langkah}
\end{column}
\end{columns}
\vspace{8pt}
\begin{columns}[T]
\begin{column}{0.47\textwidth}
\begin{Blok}{Boleh}
Membuka catatan dan modul sendiri.
\end{Blok}
\end{column}
\begin{column}{0.47\textwidth}
\begin{BlokAlert}{Belum dulu}
Bertanya ke teman atau AI. Ini simulasi kondisi ujian.
\end{BlokAlert}
\end{column}
\end{columns}
\note{Catat kesalahan yang sering muncul untuk dibahas di debrief.}
\end{frame}

\begin{frame}{Apa yang kalian temukan?}
\framesubtitle{Debrief struggle}
\begin{columns}[T]
\begin{column}{0.31\textwidth}
\begin{BlokGelap}{Pertanyaan 1}
Bagian mana yang paling lama dikerjakan?
\end{BlokGelap}
\end{column}
\begin{column}{0.31\textwidth}
\begin{BlokGelap}{Pertanyaan 2}
Bagaimana kalian membulatkan menit ke atas?
\end{BlokGelap}
\end{column}
\begin{column}{0.31\textwidth}
\begin{BlokGelap}{Pertanyaan 3}
Data uji apa yang kalian pakai?
\end{BlokGelap}
\end{column}
\end{columns}
\vspace{10pt}
Jawaban kalian akan kita uji di bagian berikutnya.\note{Tulis jawaban di papan; buktikan atau patahkan di bagian materi.}
\end{frame}

\Bagian{3}{Merangkai pertemuan 1 sampai 6}{Peta konsep, strategi soal A, dan strategi soal B}
\note{Materi 40 menit. Tunjukkan setiap contoh langsung di GDB Online.}

\begin{frame}{Peta materi UTS}
\framesubtitle{Pertemuan 1 sampai 6}
\small
\begin{tabular}{@{}>{\raggedright\arraybackslash}p{124pt}>{\raggedright\arraybackslash}p{345pt}>{\raggedright\arraybackslash}p{393pt}@{}}
\kepala{Pertemuan} & \kepala{Konsep kunci} & \kepala{Jebakan yang sering muncul}\\
P1 & struktur program, \texttt{cout}, escape & garis miring tunggal, lupa titik koma\\
\barisgenap P2 & tipe data, \texttt{cin}, \texttt{getline} & pembagian bulat, \texttt{getline} sesudah \texttt{cin}\\
P3 & \texttt{\%}, precedence, \texttt{++}, logika & urutan operasi, \texttt{x++} vs \texttt{++x}\\
\barisgenap P4 & \texttt{if}, \texttt{else if}, \texttt{switch} & \texttt{=} vs \texttt{==}, urutan syarat, \texttt{break}\\
P5 & \texttt{for}, \texttt{while}, \texttt{do-while} & akumulator tidak diinisialisasi, off-by-one\\
\barisgenap P6 & bersarang, \texttt{break}, \texttt{continue} & batas dalam, rumus pola\\
\garisdasar
\end{tabular}
\end{frame}

\begin{frame}{Strategi soal A}
\framesubtitle{Menulis program}
\begin{columns}[T]
\begin{column}{0.31\textwidth}
\begin{Langkah}{1}{Rancang dulu}
Tulis masukan, proses, keluaran, dan data uji di kertas sebelum mengetik.
\end{Langkah}
\end{column}
\begin{column}{0.31\textwidth}
\begin{Langkah}{2}{Bangun bertahap}
Baca masukan, tampilkan, kompilasi. Tambah satu bagian, kompilasi lagi.
\end{Langkah}
\end{column}
\begin{column}{0.31\textwidth}
\begin{Langkah}{3}{Uji batas}
Jalankan dengan data uji biasa dan nilai batas. Cocokkan format keluaran.
\end{Langkah}
\end{column}
\end{columns}
\vspace{8pt}
Program yang sebagian benar dan terkompilasi mendapat nilai lebih baik daripada program lengkap yang tidak terkompilasi.
\end{frame}

\begin{frame}{Strategi soal B}
\framesubtitle{Tracing dan penjelasan}
\begin{columns}[T]
\begin{column}{0.31\textwidth}
\begin{Langkah}{1}{Buat tabel}
Satu kolom per variabel, satu baris per perubahan nilai.
\end{Langkah}
\end{column}
\begin{column}{0.31\textwidth}
\begin{Langkah}{2}{Ikuti persis}
Kerjakan sesuai aturan C++: pembagian bulat, precedence, \texttt{break}.
\end{Langkah}
\end{column}
\begin{column}{0.31\textwidth}
\begin{Langkah}{3}{Jelaskan}
Pakai istilah tepat: pembagian bulat, nilai batas, akumulator, cabang.
\end{Langkah}
\end{column}
\end{columns}
\vspace{8pt}
Jawaban B tanpa langkah hanya mendapat sebagian poin, walaupun hasil akhirnya benar.
\end{frame}

\begin{frame}[fragile]{Contoh tracing gabungan}
\framesubtitle{Soal B}
\begin{columns}[T]
\begin{column}{0.55\textwidth}
\begin{Kode}[]{gabungan.cpp}
int n = 6, hasil = 0;
for (int i = 1; i <= n; i++) {
    if (i % 2 == 0) {
        hasil += i;
    } else if (i % 3 == 0) {
        hasil -= i;
    }
}
cout << hasil << endl;
\end{Kode}
\end{column}
\begin{column}{0.41\textwidth}
\only<1>{\begin{TebakDulu}
{\ITERAKickerFont\fontsize{10pt}{12pt}\selectfont\color{ITERAAlert}TEBAK DULU}\par\vspace{4pt}
Sebelum dijalankan, tulis keluarannya di kertas.
\end{TebakDulu}
}\only<2>{\KeluaranFile[]{keluaran/P07/k001.txt}
\begin{Catatan}
i = 2, 4, 6 menambah. i = 3 mengurangi. i = 6 genap, jadi cabang kedua tidak diperiksa. Hasil: 2 + 4 + 6 − 3 = 9.
\end{Catatan}
}
\end{column}
\end{columns}
\note<1>{Minta mahasiswa menebak keluaran sebelum membuka slide berikutnya. Tanyakan satu atau dua tebakan yang berbeda, lalu buktikan.}
\end{frame}

\begin{frame}[fragile]{Memecah digit}
\framesubtitle{Pola yang sering keluar}
\begin{columns}[T]
\begin{column}{0.55\textwidth}
\begin{Kode}[]{digit.cpp}
int n = 4096, jumlah = 0;
while (n > 0) {
    jumlah += n % 10;
    n /= 10;
}
cout << jumlah << endl;
\end{Kode}
\end{column}
\begin{column}{0.41\textwidth}
\only<1>{\begin{TebakDulu}
{\ITERAKickerFont\fontsize{10pt}{12pt}\selectfont\color{ITERAAlert}TEBAK DULU}\par\vspace{4pt}
Sebelum dijalankan, tulis keluarannya di kertas.
\end{TebakDulu}
}\only<2>{\KeluaranFile[]{keluaran/P07/k002.txt}
\begin{Catatan}
\texttt{\% 10} mengambil digit terakhir, \texttt{/ 10} membuangnya. Pola ini muncul di banyak soal.
\end{Catatan}
}
\end{column}
\end{columns}
\note<1>{Minta mahasiswa menebak keluaran sebelum membuka slide berikutnya. Tanyakan satu atau dua tebakan yang berbeda, lalu buktikan.}
\end{frame}

\begin{frame}[fragile]{Tebak keluarannya}
\framesubtitle{Latihan menjelaskan, 3 menit}
\begin{columns}[T]
\begin{column}{0.56\textwidth}
\begin{Kode}[]{tebak.cpp}
int a = 5, b = 0;
while (a > 0) {
    if (a % 2 == 1) {
        b += a;
    } else {
        b -= 1;
    }
    a -= 2;
}
cout << a << " " << b << endl;
\end{Kode}
\end{column}
\begin{column}{0.40\textwidth}
\begin{Catatan}
Tulis keluarannya di kertas, karakter demi karakter. Jangan dijalankan dulu.
\end{Catatan}
\end{column}
\end{columns}
\note{Contoh soal Bagian B. Beri waktu 3 menit sebelum slide jawaban.}
\end{frame}

\begin{frame}[fragile]{Tebak keluarannya: jawaban}
\framesubtitle{Latihan menjelaskan}
\begin{columns}[T]
\begin{column}{0.56\textwidth}
\begin{Kode}[]{tebak.cpp}
int a = 5, b = 0;
while (a > 0) {
    if (a % 2 == 1) {
        b += a;
    } else {
        b -= 1;
    }
    a -= 2;
}
cout << a << " " << b << endl;
\end{Kode}
\end{column}
\begin{column}{0.40\textwidth}
\begin{Keluaran}[]
-1 9
\end{Keluaran}
\begin{Catatan}
\texttt{a} bernilai 5, 3, 1, semuanya ganjil, jadi \texttt{b} = 5 + 3 + 1 = 9. Sesudah itu \texttt{a} = −1 dan perulangan berhenti.
\end{Catatan}
\end{column}
\end{columns}
\note{Tanyakan jawaban yang paling sering salah, lalu telusuri bersama.}
\end{frame}

\begin{frame}{Error yang sering muncul minggu ini}
\framesubtitle{Pesan diambil langsung dari g++}
\footnotesize
\begin{tabular}{@{}>{\raggedright\arraybackslash}p{208pt}>{\raggedright\arraybackslash}p{256pt}>{\raggedright\arraybackslash}p{398pt}@{}}
\kepala{Kesalahan} & \kepala{Contoh} & \kepala{Pesan pertama dari g++}\\
Lupa titik koma & \texttt{cout << x} & \texttt{error: expected ';' before 'return'}\\
\barisgenap Variabel belum dideklarasikan & \texttt{total += x;} & \texttt{error: 'total' was not declared in this scope}\\
= di dalam kondisi & \texttt{if (a = b)} & \texttt{warning: suggest parentheses around assignment used as truth value}\\
\barisgenap break di luar perulangan & \texttt{break; di if} & \texttt{error: break statement not within loop or switch}\\
Kurung kurawal kurang & \texttt{while \{ if \{ \}} & \texttt{error: expected '\}' at end of input}\\
\garisdasar
\end{tabular}
\vspace{10pt}

{\small Pesan di tabel ini dihasilkan dengan g++ dan opsi -Wall, sama seperti di cek.sh.}
\note{Minta mahasiswa memotret slide ini.}
\end{frame}

\Bagian{4}{Hands-on bertingkat}{45 menit, dari dasar sampai tantangan}
\note{Mahasiswa membuka folder 02 Hands-on/P7 - Review UTS - Hands-on.}

\begin{frame}{Peta hands-on}
\framesubtitle{Folder 02 Hands-on/P7 - Review UTS - Hands-on}
\small
\begin{tabular}{@{}>{\raggedright\arraybackslash}p{36pt}>{\raggedright\arraybackslash}p{267pt}>{\raggedright\arraybackslash}p{110pt}>{\raggedright\arraybackslash}p{83pt}>{\raggedright\arraybackslash}p{341pt}@{}}
\kepala{No} & \kepala{File} & \kepala{Tingkat} & \kepala{Waktu} & \kepala{Yang dilatih}\\
1 & \texttt{handson1-parkir.cpp} & Dasar & 20' & masukan, \texttt{/} dan \texttt{\%}, \texttt{else if}\\
\barisgenap 2 & \texttt{handson2-digit.cpp} & Menengah & 15' & \texttt{while} dengan pola digit\\
3 & \texttt{handson3-telusur.cpp} & Tantangan & 10' & trace table bersarang dengan \texttt{continue} dan \texttt{break}\\
\barisgenap + & \texttt{bonus-menu.cpp} & Pengayaan & sisa & menu \texttt{do-while}, validasi\\
\garisdasar
\end{tabular}
\vspace{10pt}

Kerjakan berurutan. Cocokkan keluaran dengan folder \texttt{expected/}; latihan yang membaca masukan memakai data di folder \texttt{input/}.\note{Saat berkeliling, minta mahasiswa yang mengerjakan latihan tantangan menjelaskan blok PENJELASAN mereka.}
\end{frame}

\begin{frame}{Cara mengecek dan aturannya}
\framesubtitle{Hands-on}
\begin{columns}[T]
\begin{column}{0.47\textwidth}
\begin{Blok}{Di GDB Online}
Salin isi file latihan, lengkapi TODO, klik Run. Bila program membaca masukan, ketik isi file di \texttt{input/}. Bandingkan dengan \texttt{expected/}.
\end{Blok}
\end{column}
\begin{column}{0.47\textwidth}
\begin{BlokGelap}{Di komputer sendiri}
Jalankan \texttt{bash cek.sh}. Skrip mengompilasi dengan \texttt{-Wall -Werror}, memberi masukan dari \texttt{input/}, lalu mencocokkan keluaran.
\end{BlokGelap}
\end{column}
\end{columns}
\vspace{6pt}
\begin{Catatan}
AI boleh untuk memahami pesan error, tidak untuk meminta solusi utuh. Exit Ticket dikerjakan tanpa bantuan apa pun.
\end{Catatan}
\note{Hands-on tidak dinilai langsung; yang dinilai Exit Ticket dan tugas.}
\end{frame}

\Bagian{5}{Exit Ticket dan tugas}{25 menit terakhir, lalu pekerjaan rumah}
\note{Mulai Exit Ticket tepat waktu.}

\begin{frame}{Exit Ticket 7}
\framesubtitle{Individual, 25 menit, tanpa AI dan internet}
\small
\begin{tabular}{@{}>{\raggedright\arraybackslash}p{60pt}>{\raggedright\arraybackslash}p{560pt}>{\raggedright\arraybackslash}p{80pt}>{\raggedright\arraybackslash}p{150pt}@{}}
\kepala{Soal} & \kepala{Isi} & \kepala{Poin} & \kepala{CPMK}\\
A1 & Tulis program singkat yang memakai masukan, percabangan, dan perulangan & 25 & CPMK0607\\
\barisgenap A2 & Perbaiki program agar keluarannya sesuai contoh & 25 & CPMK0607\\
B1 & Buat trace table untuk perulangan bersarang dengan \texttt{continue} & 20 & CPMK0608\\
\barisgenap B2 & Jelaskan mengapa dua cabang tidak dijalankan bersamaan pada sebuah rantai \texttt{else if} & 20 & CPMK0608\\
B3 & Sebutkan materi pertemuan yang dipakai sebuah soal dan alasannya & 10 & CPMK0608\\
\garisdasar
\end{tabular}
\vspace{10pt}

{\small Soal A dinilai dari keluaran yang tepat sama. Soal B dinilai dari ketepatan dan alasan. Pelanggaran aturan membuat nilai Exit Ticket ini 0.}
\note{Soal lengkap dibagikan terpisah dan tidak disimpan di repo publik.}
\end{frame}

\begin{frame}{Sebelum Pertemuan 8}
\framesubtitle{Persiapan}
\begin{columns}[T]
\begin{column}{0.31\textwidth}
\begin{Langkah}{1}{Selesaikan}
Hands-on yang belum lulus \texttt{cek.sh}, terutama latihan tantangan.
\end{Langkah}
\end{column}
\begin{column}{0.31\textwidth}
\begin{Langkah}{2}{Latih}
Tidak ada tugas minggu ini. Kerjakan ulang ketiga simulasi dengan batas waktu ujian.
\end{Langkah}
\end{column}
\begin{column}{0.31\textwidth}
\begin{Langkah}{3}{Baca}
Modul Belajar 7 untuk mengulang, lalu bagian awal modul berikutnya.
\end{Langkah}
\end{column}
\end{columns}
\vspace{10pt}
Pertemuan 8 adalah UTS. Kisi-kisi, tata tertib, dan contoh soal ada di slide P8 dan Modul Belajar 8.\note{Tanyakan satu hal yang paling membingungkan hari ini untuk dibuka di review minggu depan.}
\end{frame}

\SlidePenutup{Terima kasih}{Sampai jumpa di UTS. Kerjakan contoh soal P8 dengan batas waktu.}
\note{Slide penutup.}
\end{document}
"
\documentclass[t]{beamer}
\usetheme{ITERA}
% Mode catatan pembicara: aktif bila \catatan didefinisikan sebelum \input.
\ifdefined\catatan
  \setbeameroption{show only notes}
  \setbeamertemplate{note page}{\vspace*{20pt}\hspace*{4pt}\begin{minipage}{900pt}
    {\ITERAKickerFont\fontsize{11pt}{14pt}\selectfont\color{ITERAGoldText}CATATAN PEMBICARA \ \ |\ \ SLIDE \insertframenumber}\par\vspace{4pt}
    {\ITERADisplay\bfseries\fontsize{22pt}{28pt}\selectfont\color{ITERAStructure}\insertframetitle}\par\vspace{12pt}
    \fontsize{15pt}{23pt}\selectfont\insertnote\end{minipage}}
\fi
\tikzset{fase/.style={draw=none,fill=#1,rounded corners=3pt,minimum width=158pt,minimum height=118pt,text width=146pt,align=center}}

\renewcommand\ITERAmeeting{Pertemuan 7}
\begin{document}
\SlideJudul{IF25-11002 PRAKTIKUM PEMROGRAMAN}{Pertemuan 7\\Review dan Simulasi UTS}{Merangkai materi Pertemuan 1 sampai 6 dan berlatih dengan soal bergaya UTS}{Tim Pengajar}{Tahun 2026. Sejajar dengan Pertemuan 7 Algoritma Pemrograman: Review dan Simulasi UTS.}
\note{Buka dengan menanyakan satu hal dari pertemuan lalu yang masih membingungkan.}

\begin{frame}{Dua mata kuliah, satu minggu}
\framesubtitle{Sebelum mulai}
Topik minggu ini sama di dua mata kuliah, dengan peran yang berbeda.
\vspace{10pt}

\begin{columns}[T]
\begin{column}{0.47\textwidth}
\begin{Blok}{Algoritma: Review dan Simulasi UTS}
Review pseudocode, flowchart, dan trace table untuk materi pertemuan 1 sampai 6.
\end{Blok}
\end{column}
\begin{column}{0.47\textwidth}
\begin{BlokGelap}{Praktikum: Review UTS}
Review C++ pertemuan 1 sampai 6 dan simulasi soal A (kode) dan soal B (tracing dan penjelasan).
\end{BlokGelap}
\end{column}
\end{columns}
\note{Tanyakan apa yang sudah dibahas di Algoritma minggu ini. Gunakan jawabannya sebagai jembatan ke praktikum.}
\end{frame}

\begin{frame}{Saat pulang nanti, kalian bisa}
\framesubtitle{Target hari ini}
\begin{columns}[T]
\begin{column}{0.47\textwidth}
\begin{Langkah}{1}{Menerapkan}
Menyelesaikan masalah yang menggabungkan masukan, ekspresi, percabangan, dan perulangan dalam satu program C++ dalam batas waktu ujian

{\footnotesize\color{ITERAInkMuted}Sub-CPMK 7.1, CPMK0607}
\end{Langkah}
\end{column}
\begin{column}{0.47\textwidth}
\begin{Langkah}{2}{Menjelaskan}
Menelusuri program yang menggabungkan percabangan dan perulangan bersarang, lalu menjelaskan setiap perubahan nilai dengan istilah yang tepat

{\footnotesize\color{ITERAInkMuted}Sub-CPMK 7.2, CPMK0608}
\end{Langkah}
\end{column}
\end{columns}
\vspace{8pt}
\begin{Catatan}
Dua kemampuan ini dinilai sama berat di Exit Ticket, tugas, UTS, dan UAS.
\end{Catatan}
\note{Bacakan target dengan pelan. Kata kuncinya menerapkan dan menjelaskan.}
\end{frame}

\begin{frame}{Ingat minggu lalu}
\framesubtitle{Review, 5 menit}
\begin{columns}[T]
\begin{column}{0.31\textwidth}
\begin{BlokGelap}{Pertanyaan 1}
Apa rumus banyak bintang pada piramida baris b?
\end{BlokGelap}
\end{column}
\begin{column}{0.31\textwidth}
\begin{BlokGelap}{Pertanyaan 2}
Apa beda \texttt{break} dan \texttt{continue} di perulangan bersarang?
\end{BlokGelap}
\end{column}
\begin{column}{0.31\textwidth}
\begin{BlokGelap}{Pertanyaan 3}
Bagaimana menghitung jam dari menit dengan pembulatan ke atas?
\end{BlokGelap}
\end{column}
\end{columns}
\vspace{8pt}
Jawab di kertas dulu, lalu cocokkan dengan teman sebelah.\note{Gunakan tiga pertanyaan ini sebagai diagnosis cepat. Catat topik yang paling banyak salah untuk dibahas di bagian materi.}
\end{frame}

\Bagian{1}{Masalah pembuka}{Satu soal, semua materi}
\note{Masalah pembuka 10 menit.}

\begin{frame}[fragile]{Satu soal, semua materi}
\framesubtitle{Masalah minggu ini}
\begin{columns}[T]
\begin{column}{0.55\textwidth}
{\small Soal UTS jarang menguji satu topik saja. Contohnya tarif parkir: membaca masukan, membulatkan menit ke jam dengan \texttt{/} dan \texttt{\%}, memilih tarif dengan \texttt{if}, dan mengulang untuk beberapa kendaraan.

\vspace{6pt}
Hari ini kita berlatih memecah soal seperti itu menjadi bagian-bagian yang sudah kalian kuasai.\par}
\vspace{6pt}
\begin{Catatan}
\textbf{Diskusi 2 menit.} Materi pertemuan mana saja yang dipakai soal ini? Bagian mana yang menurut kalian paling rawan salah?
\end{Catatan}
\end{column}
\begin{column}{0.41\textwidth}
\begin{Keluaran}[]
M 61   -> 2 jam, tarif 3000
B 180  -> 3 jam, tarif 11000
X 10   -> Jenis tidak dikenal
\end{Keluaran}
\end{column}
\end{columns}
\note{Tampung beberapa jawaban tanpa menilai benar salah.}
\end{frame}

\begin{frame}{Dari langkah ke instruksi}
\framesubtitle{Sambungan dengan Algoritma}
\begin{columns}[T]
\begin{column}{0.31\textwidth}
\begin{Langkah}{1}{Baca soal}
Garis bawahi masukan, keluaran, dan aturan.
\end{Langkah}
\end{column}
\begin{column}{0.31\textwidth}
\begin{Langkah}{2}{Pecah}
Petakan setiap aturan ke konstruksi C++ yang sudah dipelajari.
\end{Langkah}
\end{column}
\begin{column}{0.31\textwidth}
\begin{Langkah}{3}{Uji}
Siapkan data uji, termasuk nilai batas, sebelum menulis kode.
\end{Langkah}
\end{column}
\end{columns}
\vspace{10pt}
Langkah ini sama dengan yang dilatih di Algoritma minggu ini, hanya hasil akhirnya kode C++.\note{Hubungkan eksplisit ke Algoritma minggu ini.}
\end{frame}

\Bagian{2}{Coba dulu, baru dijelaskan}{25 menit eksperimen di GDB Online}
\note{Struggle 25 menit. Berkeliling, jangan menjelaskan materi dulu.}

\begin{frame}{Misi kalian}
\framesubtitle{Struggle, 25 menit}
\begin{columns}[T]
\begin{column}{0.31\textwidth}
\begin{Langkah}{1}{Pilih}
Kerjakan \texttt{handson1-parkir.cpp} tanpa bantuan selama 20 menit.
\end{Langkah}
\end{column}
\begin{column}{0.31\textwidth}
\begin{Langkah}{2}{Uji}
Uji dengan 60 dan 61 menit. Apakah hasilnya berbeda satu jam?
\end{Langkah}
\end{column}
\begin{column}{0.31\textwidth}
\begin{Langkah}{3}{Catat}
Tulis satu hal yang membuat kalian tersendat.
\end{Langkah}
\end{column}
\end{columns}
\vspace{8pt}
\begin{columns}[T]
\begin{column}{0.47\textwidth}
\begin{Blok}{Boleh}
Membuka catatan dan modul sendiri.
\end{Blok}
\end{column}
\begin{column}{0.47\textwidth}
\begin{BlokAlert}{Belum dulu}
Bertanya ke teman atau AI. Ini simulasi kondisi ujian.
\end{BlokAlert}
\end{column}
\end{columns}
\note{Catat kesalahan yang sering muncul untuk dibahas di debrief.}
\end{frame}

\begin{frame}{Apa yang kalian temukan?}
\framesubtitle{Debrief struggle}
\begin{columns}[T]
\begin{column}{0.31\textwidth}
\begin{BlokGelap}{Pertanyaan 1}
Bagian mana yang paling lama dikerjakan?
\end{BlokGelap}
\end{column}
\begin{column}{0.31\textwidth}
\begin{BlokGelap}{Pertanyaan 2}
Bagaimana kalian membulatkan menit ke atas?
\end{BlokGelap}
\end{column}
\begin{column}{0.31\textwidth}
\begin{BlokGelap}{Pertanyaan 3}
Data uji apa yang kalian pakai?
\end{BlokGelap}
\end{column}
\end{columns}
\vspace{10pt}
Jawaban kalian akan kita uji di bagian berikutnya.\note{Tulis jawaban di papan; buktikan atau patahkan di bagian materi.}
\end{frame}

\Bagian{3}{Merangkai pertemuan 1 sampai 6}{Peta konsep, strategi soal A, dan strategi soal B}
\note{Materi 40 menit. Tunjukkan setiap contoh langsung di GDB Online.}

\begin{frame}{Peta materi UTS}
\framesubtitle{Pertemuan 1 sampai 6}
\small
\begin{tabular}{@{}>{\raggedright\arraybackslash}p{124pt}>{\raggedright\arraybackslash}p{345pt}>{\raggedright\arraybackslash}p{393pt}@{}}
\kepala{Pertemuan} & \kepala{Konsep kunci} & \kepala{Jebakan yang sering muncul}\\
P1 & struktur program, \texttt{cout}, escape & garis miring tunggal, lupa titik koma\\
\barisgenap P2 & tipe data, \texttt{cin}, \texttt{getline} & pembagian bulat, \texttt{getline} sesudah \texttt{cin}\\
P3 & \texttt{\%}, precedence, \texttt{++}, logika & urutan operasi, \texttt{x++} vs \texttt{++x}\\
\barisgenap P4 & \texttt{if}, \texttt{else if}, \texttt{switch} & \texttt{=} vs \texttt{==}, urutan syarat, \texttt{break}\\
P5 & \texttt{for}, \texttt{while}, \texttt{do-while} & akumulator tidak diinisialisasi, off-by-one\\
\barisgenap P6 & bersarang, \texttt{break}, \texttt{continue} & batas dalam, rumus pola\\
\garisdasar
\end{tabular}
\end{frame}

\begin{frame}{Strategi soal A}
\framesubtitle{Menulis program}
\begin{columns}[T]
\begin{column}{0.31\textwidth}
\begin{Langkah}{1}{Rancang dulu}
Tulis masukan, proses, keluaran, dan data uji di kertas sebelum mengetik.
\end{Langkah}
\end{column}
\begin{column}{0.31\textwidth}
\begin{Langkah}{2}{Bangun bertahap}
Baca masukan, tampilkan, kompilasi. Tambah satu bagian, kompilasi lagi.
\end{Langkah}
\end{column}
\begin{column}{0.31\textwidth}
\begin{Langkah}{3}{Uji batas}
Jalankan dengan data uji biasa dan nilai batas. Cocokkan format keluaran.
\end{Langkah}
\end{column}
\end{columns}
\vspace{8pt}
Program yang sebagian benar dan terkompilasi mendapat nilai lebih baik daripada program lengkap yang tidak terkompilasi.
\end{frame}

\begin{frame}{Strategi soal B}
\framesubtitle{Tracing dan penjelasan}
\begin{columns}[T]
\begin{column}{0.31\textwidth}
\begin{Langkah}{1}{Buat tabel}
Satu kolom per variabel, satu baris per perubahan nilai.
\end{Langkah}
\end{column}
\begin{column}{0.31\textwidth}
\begin{Langkah}{2}{Ikuti persis}
Kerjakan sesuai aturan C++: pembagian bulat, precedence, \texttt{break}.
\end{Langkah}
\end{column}
\begin{column}{0.31\textwidth}
\begin{Langkah}{3}{Jelaskan}
Pakai istilah tepat: pembagian bulat, nilai batas, akumulator, cabang.
\end{Langkah}
\end{column}
\end{columns}
\vspace{8pt}
Jawaban B tanpa langkah hanya mendapat sebagian poin, walaupun hasil akhirnya benar.
\end{frame}

\begin{frame}[fragile]{Contoh tracing gabungan}
\framesubtitle{Soal B}
\begin{columns}[T]
\begin{column}{0.55\textwidth}
\begin{Kode}[]{gabungan.cpp}
int n = 6, hasil = 0;
for (int i = 1; i <= n; i++) {
    if (i % 2 == 0) {
        hasil += i;
    } else if (i % 3 == 0) {
        hasil -= i;
    }
}
cout << hasil << endl;
\end{Kode}
\end{column}
\begin{column}{0.41\textwidth}
\only<1>{\begin{TebakDulu}
{\ITERAKickerFont\fontsize{10pt}{12pt}\selectfont\color{ITERAAlert}TEBAK DULU}\par\vspace{4pt}
Sebelum dijalankan, tulis keluarannya di kertas.
\end{TebakDulu}
}\only<2>{\KeluaranFile[]{keluaran/P07/k001.txt}
\begin{Catatan}
i = 2, 4, 6 menambah. i = 3 mengurangi. i = 6 genap, jadi cabang kedua tidak diperiksa. Hasil: 2 + 4 + 6 − 3 = 9.
\end{Catatan}
}
\end{column}
\end{columns}
\note<1>{Minta mahasiswa menebak keluaran sebelum membuka slide berikutnya. Tanyakan satu atau dua tebakan yang berbeda, lalu buktikan.}
\end{frame}

\begin{frame}[fragile]{Memecah digit}
\framesubtitle{Pola yang sering keluar}
\begin{columns}[T]
\begin{column}{0.55\textwidth}
\begin{Kode}[]{digit.cpp}
int n = 4096, jumlah = 0;
while (n > 0) {
    jumlah += n % 10;
    n /= 10;
}
cout << jumlah << endl;
\end{Kode}
\end{column}
\begin{column}{0.41\textwidth}
\only<1>{\begin{TebakDulu}
{\ITERAKickerFont\fontsize{10pt}{12pt}\selectfont\color{ITERAAlert}TEBAK DULU}\par\vspace{4pt}
Sebelum dijalankan, tulis keluarannya di kertas.
\end{TebakDulu}
}\only<2>{\KeluaranFile[]{keluaran/P07/k002.txt}
\begin{Catatan}
\texttt{\% 10} mengambil digit terakhir, \texttt{/ 10} membuangnya. Pola ini muncul di banyak soal.
\end{Catatan}
}
\end{column}
\end{columns}
\note<1>{Minta mahasiswa menebak keluaran sebelum membuka slide berikutnya. Tanyakan satu atau dua tebakan yang berbeda, lalu buktikan.}
\end{frame}

\begin{frame}[fragile]{Tebak keluarannya}
\framesubtitle{Latihan menjelaskan, 3 menit}
\begin{columns}[T]
\begin{column}{0.56\textwidth}
\begin{Kode}[]{tebak.cpp}
int a = 5, b = 0;
while (a > 0) {
    if (a % 2 == 1) {
        b += a;
    } else {
        b -= 1;
    }
    a -= 2;
}
cout << a << " " << b << endl;
\end{Kode}
\end{column}
\begin{column}{0.40\textwidth}
\begin{Catatan}
Tulis keluarannya di kertas, karakter demi karakter. Jangan dijalankan dulu.
\end{Catatan}
\end{column}
\end{columns}
\note{Contoh soal Bagian B. Beri waktu 3 menit sebelum slide jawaban.}
\end{frame}

\begin{frame}[fragile]{Tebak keluarannya: jawaban}
\framesubtitle{Latihan menjelaskan}
\begin{columns}[T]
\begin{column}{0.56\textwidth}
\begin{Kode}[]{tebak.cpp}
int a = 5, b = 0;
while (a > 0) {
    if (a % 2 == 1) {
        b += a;
    } else {
        b -= 1;
    }
    a -= 2;
}
cout << a << " " << b << endl;
\end{Kode}
\end{column}
\begin{column}{0.40\textwidth}
\begin{Keluaran}[]
-1 9
\end{Keluaran}
\begin{Catatan}
\texttt{a} bernilai 5, 3, 1, semuanya ganjil, jadi \texttt{b} = 5 + 3 + 1 = 9. Sesudah itu \texttt{a} = −1 dan perulangan berhenti.
\end{Catatan}
\end{column}
\end{columns}
\note{Tanyakan jawaban yang paling sering salah, lalu telusuri bersama.}
\end{frame}

\begin{frame}{Error yang sering muncul minggu ini}
\framesubtitle{Pesan diambil langsung dari g++}
\footnotesize
\begin{tabular}{@{}>{\raggedright\arraybackslash}p{208pt}>{\raggedright\arraybackslash}p{256pt}>{\raggedright\arraybackslash}p{398pt}@{}}
\kepala{Kesalahan} & \kepala{Contoh} & \kepala{Pesan pertama dari g++}\\
Lupa titik koma & \texttt{cout << x} & \texttt{error: expected ';' before 'return'}\\
\barisgenap Variabel belum dideklarasikan & \texttt{total += x;} & \texttt{error: 'total' was not declared in this scope}\\
= di dalam kondisi & \texttt{if (a = b)} & \texttt{warning: suggest parentheses around assignment used as truth value}\\
\barisgenap break di luar perulangan & \texttt{break; di if} & \texttt{error: break statement not within loop or switch}\\
Kurung kurawal kurang & \texttt{while \{ if \{ \}} & \texttt{error: expected '\}' at end of input}\\
\garisdasar
\end{tabular}
\vspace{10pt}

{\small Pesan di tabel ini dihasilkan dengan g++ dan opsi -Wall, sama seperti di cek.sh.}
\note{Minta mahasiswa memotret slide ini.}
\end{frame}

\Bagian{4}{Hands-on bertingkat}{45 menit, dari dasar sampai tantangan}
\note{Mahasiswa membuka folder 02 Hands-on/P7 - Review UTS - Hands-on.}

\begin{frame}{Peta hands-on}
\framesubtitle{Folder 02 Hands-on/P7 - Review UTS - Hands-on}
\small
\begin{tabular}{@{}>{\raggedright\arraybackslash}p{36pt}>{\raggedright\arraybackslash}p{267pt}>{\raggedright\arraybackslash}p{110pt}>{\raggedright\arraybackslash}p{83pt}>{\raggedright\arraybackslash}p{341pt}@{}}
\kepala{No} & \kepala{File} & \kepala{Tingkat} & \kepala{Waktu} & \kepala{Yang dilatih}\\
1 & \texttt{handson1-parkir.cpp} & Dasar & 20' & masukan, \texttt{/} dan \texttt{\%}, \texttt{else if}\\
\barisgenap 2 & \texttt{handson2-digit.cpp} & Menengah & 15' & \texttt{while} dengan pola digit\\
3 & \texttt{handson3-telusur.cpp} & Tantangan & 10' & trace table bersarang dengan \texttt{continue} dan \texttt{break}\\
\barisgenap + & \texttt{bonus-menu.cpp} & Pengayaan & sisa & menu \texttt{do-while}, validasi\\
\garisdasar
\end{tabular}
\vspace{10pt}

Kerjakan berurutan. Cocokkan keluaran dengan folder \texttt{expected/}; latihan yang membaca masukan memakai data di folder \texttt{input/}.\note{Saat berkeliling, minta mahasiswa yang mengerjakan latihan tantangan menjelaskan blok PENJELASAN mereka.}
\end{frame}

\begin{frame}{Cara mengecek dan aturannya}
\framesubtitle{Hands-on}
\begin{columns}[T]
\begin{column}{0.47\textwidth}
\begin{Blok}{Di GDB Online}
Salin isi file latihan, lengkapi TODO, klik Run. Bila program membaca masukan, ketik isi file di \texttt{input/}. Bandingkan dengan \texttt{expected/}.
\end{Blok}
\end{column}
\begin{column}{0.47\textwidth}
\begin{BlokGelap}{Di komputer sendiri}
Jalankan \texttt{bash cek.sh}. Skrip mengompilasi dengan \texttt{-Wall -Werror}, memberi masukan dari \texttt{input/}, lalu mencocokkan keluaran.
\end{BlokGelap}
\end{column}
\end{columns}
\vspace{6pt}
\begin{Catatan}
AI boleh untuk memahami pesan error, tidak untuk meminta solusi utuh. Exit Ticket dikerjakan tanpa bantuan apa pun.
\end{Catatan}
\note{Hands-on tidak dinilai langsung; yang dinilai Exit Ticket dan tugas.}
\end{frame}

\Bagian{5}{Exit Ticket dan tugas}{25 menit terakhir, lalu pekerjaan rumah}
\note{Mulai Exit Ticket tepat waktu.}

\begin{frame}{Exit Ticket 7}
\framesubtitle{Individual, 25 menit, tanpa AI dan internet}
\small
\begin{tabular}{@{}>{\raggedright\arraybackslash}p{60pt}>{\raggedright\arraybackslash}p{560pt}>{\raggedright\arraybackslash}p{80pt}>{\raggedright\arraybackslash}p{150pt}@{}}
\kepala{Soal} & \kepala{Isi} & \kepala{Poin} & \kepala{CPMK}\\
A1 & Tulis program singkat yang memakai masukan, percabangan, dan perulangan & 25 & CPMK0607\\
\barisgenap A2 & Perbaiki program agar keluarannya sesuai contoh & 25 & CPMK0607\\
B1 & Buat trace table untuk perulangan bersarang dengan \texttt{continue} & 20 & CPMK0608\\
\barisgenap B2 & Jelaskan mengapa dua cabang tidak dijalankan bersamaan pada sebuah rantai \texttt{else if} & 20 & CPMK0608\\
B3 & Sebutkan materi pertemuan yang dipakai sebuah soal dan alasannya & 10 & CPMK0608\\
\garisdasar
\end{tabular}
\vspace{10pt}

{\small Soal A dinilai dari keluaran yang tepat sama. Soal B dinilai dari ketepatan dan alasan. Pelanggaran aturan membuat nilai Exit Ticket ini 0.}
\note{Soal lengkap dibagikan terpisah dan tidak disimpan di repo publik.}
\end{frame}

\begin{frame}{Sebelum Pertemuan 8}
\framesubtitle{Persiapan}
\begin{columns}[T]
\begin{column}{0.31\textwidth}
\begin{Langkah}{1}{Selesaikan}
Hands-on yang belum lulus \texttt{cek.sh}, terutama latihan tantangan.
\end{Langkah}
\end{column}
\begin{column}{0.31\textwidth}
\begin{Langkah}{2}{Latih}
Tidak ada tugas minggu ini. Kerjakan ulang ketiga simulasi dengan batas waktu ujian.
\end{Langkah}
\end{column}
\begin{column}{0.31\textwidth}
\begin{Langkah}{3}{Baca}
Modul Belajar 7 untuk mengulang, lalu bagian awal modul berikutnya.
\end{Langkah}
\end{column}
\end{columns}
\vspace{10pt}
Pertemuan 8 adalah UTS. Kisi-kisi, tata tertib, dan contoh soal ada di slide P8 dan Modul Belajar 8.\note{Tanyakan satu hal yang paling membingungkan hari ini untuk dibuka di review minggu depan.}
\end{frame}

\SlidePenutup{Terima kasih}{Sampai jumpa di UTS. Kerjakan contoh soal P8 dengan batas waktu.}
\note{Slide penutup.}
\end{document}
"
\documentclass[t]{beamer}
\usetheme{ITERA}
% Mode catatan pembicara: aktif bila \catatan didefinisikan sebelum \input.
\ifdefined\catatan
  \setbeameroption{show only notes}
  \setbeamertemplate{note page}{\vspace*{20pt}\hspace*{4pt}\begin{minipage}{900pt}
    {\ITERAKickerFont\fontsize{11pt}{14pt}\selectfont\color{ITERAGoldText}CATATAN PEMBICARA \ \ |\ \ SLIDE \insertframenumber}\par\vspace{4pt}
    {\ITERADisplay\bfseries\fontsize{22pt}{28pt}\selectfont\color{ITERAStructure}\insertframetitle}\par\vspace{12pt}
    \fontsize{15pt}{23pt}\selectfont\insertnote\end{minipage}}
\fi
\tikzset{fase/.style={draw=none,fill=#1,rounded corners=3pt,minimum width=158pt,minimum height=118pt,text width=146pt,align=center}}

\renewcommand\ITERAmeeting{Pertemuan 7}
\begin{document}
\SlideJudul{IF25-11002 PRAKTIKUM PEMROGRAMAN}{Pertemuan 7\\Review dan Simulasi UTS}{Merangkai materi Pertemuan 1 sampai 6 dan berlatih dengan soal bergaya UTS}{Tim Pengajar}{Tahun 2026. Sejajar dengan Pertemuan 7 Algoritma Pemrograman: Review dan Simulasi UTS.}
\note{Buka dengan menanyakan satu hal dari pertemuan lalu yang masih membingungkan.}

\begin{frame}{Dua mata kuliah, satu minggu}
\framesubtitle{Sebelum mulai}
Topik minggu ini sama di dua mata kuliah, dengan peran yang berbeda.
\vspace{10pt}

\begin{columns}[T]
\begin{column}{0.47\textwidth}
\begin{Blok}{Algoritma: Review dan Simulasi UTS}
Review pseudocode, flowchart, dan trace table untuk materi pertemuan 1 sampai 6.
\end{Blok}
\end{column}
\begin{column}{0.47\textwidth}
\begin{BlokGelap}{Praktikum: Review UTS}
Review C++ pertemuan 1 sampai 6 dan simulasi soal A (kode) dan soal B (tracing dan penjelasan).
\end{BlokGelap}
\end{column}
\end{columns}
\note{Tanyakan apa yang sudah dibahas di Algoritma minggu ini. Gunakan jawabannya sebagai jembatan ke praktikum.}
\end{frame}

\begin{frame}{Saat pulang nanti, kalian bisa}
\framesubtitle{Target hari ini}
\begin{columns}[T]
\begin{column}{0.47\textwidth}
\begin{Langkah}{1}{Menerapkan}
Menyelesaikan masalah yang menggabungkan masukan, ekspresi, percabangan, dan perulangan dalam satu program C++ dalam batas waktu ujian

{\footnotesize\color{ITERAInkMuted}Sub-CPMK 7.1, CPMK0607}
\end{Langkah}
\end{column}
\begin{column}{0.47\textwidth}
\begin{Langkah}{2}{Menjelaskan}
Menelusuri program yang menggabungkan percabangan dan perulangan bersarang, lalu menjelaskan setiap perubahan nilai dengan istilah yang tepat

{\footnotesize\color{ITERAInkMuted}Sub-CPMK 7.2, CPMK0608}
\end{Langkah}
\end{column}
\end{columns}
\vspace{8pt}
\begin{Catatan}
Dua kemampuan ini dinilai sama berat di Exit Ticket, tugas, UTS, dan UAS.
\end{Catatan}
\note{Bacakan target dengan pelan. Kata kuncinya menerapkan dan menjelaskan.}
\end{frame}

\begin{frame}{Ingat minggu lalu}
\framesubtitle{Review, 5 menit}
\begin{columns}[T]
\begin{column}{0.31\textwidth}
\begin{BlokGelap}{Pertanyaan 1}
Apa rumus banyak bintang pada piramida baris b?
\end{BlokGelap}
\end{column}
\begin{column}{0.31\textwidth}
\begin{BlokGelap}{Pertanyaan 2}
Apa beda \texttt{break} dan \texttt{continue} di perulangan bersarang?
\end{BlokGelap}
\end{column}
\begin{column}{0.31\textwidth}
\begin{BlokGelap}{Pertanyaan 3}
Bagaimana menghitung jam dari menit dengan pembulatan ke atas?
\end{BlokGelap}
\end{column}
\end{columns}
\vspace{8pt}
Jawab di kertas dulu, lalu cocokkan dengan teman sebelah.\note{Gunakan tiga pertanyaan ini sebagai diagnosis cepat. Catat topik yang paling banyak salah untuk dibahas di bagian materi.}
\end{frame}

\Bagian{1}{Masalah pembuka}{Satu soal, semua materi}
\note{Masalah pembuka 10 menit.}

\begin{frame}[fragile]{Satu soal, semua materi}
\framesubtitle{Masalah minggu ini}
\begin{columns}[T]
\begin{column}{0.55\textwidth}
{\small Soal UTS jarang menguji satu topik saja. Contohnya tarif parkir: membaca masukan, membulatkan menit ke jam dengan \texttt{/} dan \texttt{\%}, memilih tarif dengan \texttt{if}, dan mengulang untuk beberapa kendaraan.

\vspace{6pt}
Hari ini kita berlatih memecah soal seperti itu menjadi bagian-bagian yang sudah kalian kuasai.\par}
\vspace{6pt}
\begin{Catatan}
\textbf{Diskusi 2 menit.} Materi pertemuan mana saja yang dipakai soal ini? Bagian mana yang menurut kalian paling rawan salah?
\end{Catatan}
\end{column}
\begin{column}{0.41\textwidth}
\begin{Keluaran}[]
M 61   -> 2 jam, tarif 3000
B 180  -> 3 jam, tarif 11000
X 10   -> Jenis tidak dikenal
\end{Keluaran}
\end{column}
\end{columns}
\note{Tampung beberapa jawaban tanpa menilai benar salah.}
\end{frame}

\begin{frame}{Dari langkah ke instruksi}
\framesubtitle{Sambungan dengan Algoritma}
\begin{columns}[T]
\begin{column}{0.31\textwidth}
\begin{Langkah}{1}{Baca soal}
Garis bawahi masukan, keluaran, dan aturan.
\end{Langkah}
\end{column}
\begin{column}{0.31\textwidth}
\begin{Langkah}{2}{Pecah}
Petakan setiap aturan ke konstruksi C++ yang sudah dipelajari.
\end{Langkah}
\end{column}
\begin{column}{0.31\textwidth}
\begin{Langkah}{3}{Uji}
Siapkan data uji, termasuk nilai batas, sebelum menulis kode.
\end{Langkah}
\end{column}
\end{columns}
\vspace{10pt}
Langkah ini sama dengan yang dilatih di Algoritma minggu ini, hanya hasil akhirnya kode C++.\note{Hubungkan eksplisit ke Algoritma minggu ini.}
\end{frame}

\Bagian{2}{Coba dulu, baru dijelaskan}{25 menit eksperimen di GDB Online}
\note{Struggle 25 menit. Berkeliling, jangan menjelaskan materi dulu.}

\begin{frame}{Misi kalian}
\framesubtitle{Struggle, 25 menit}
\begin{columns}[T]
\begin{column}{0.31\textwidth}
\begin{Langkah}{1}{Pilih}
Kerjakan \texttt{handson1-parkir.cpp} tanpa bantuan selama 20 menit.
\end{Langkah}
\end{column}
\begin{column}{0.31\textwidth}
\begin{Langkah}{2}{Uji}
Uji dengan 60 dan 61 menit. Apakah hasilnya berbeda satu jam?
\end{Langkah}
\end{column}
\begin{column}{0.31\textwidth}
\begin{Langkah}{3}{Catat}
Tulis satu hal yang membuat kalian tersendat.
\end{Langkah}
\end{column}
\end{columns}
\vspace{8pt}
\begin{columns}[T]
\begin{column}{0.47\textwidth}
\begin{Blok}{Boleh}
Membuka catatan dan modul sendiri.
\end{Blok}
\end{column}
\begin{column}{0.47\textwidth}
\begin{BlokAlert}{Belum dulu}
Bertanya ke teman atau AI. Ini simulasi kondisi ujian.
\end{BlokAlert}
\end{column}
\end{columns}
\note{Catat kesalahan yang sering muncul untuk dibahas di debrief.}
\end{frame}

\begin{frame}{Apa yang kalian temukan?}
\framesubtitle{Debrief struggle}
\begin{columns}[T]
\begin{column}{0.31\textwidth}
\begin{BlokGelap}{Pertanyaan 1}
Bagian mana yang paling lama dikerjakan?
\end{BlokGelap}
\end{column}
\begin{column}{0.31\textwidth}
\begin{BlokGelap}{Pertanyaan 2}
Bagaimana kalian membulatkan menit ke atas?
\end{BlokGelap}
\end{column}
\begin{column}{0.31\textwidth}
\begin{BlokGelap}{Pertanyaan 3}
Data uji apa yang kalian pakai?
\end{BlokGelap}
\end{column}
\end{columns}
\vspace{10pt}
Jawaban kalian akan kita uji di bagian berikutnya.\note{Tulis jawaban di papan; buktikan atau patahkan di bagian materi.}
\end{frame}

\Bagian{3}{Merangkai pertemuan 1 sampai 6}{Peta konsep, strategi soal A, dan strategi soal B}
\note{Materi 40 menit. Tunjukkan setiap contoh langsung di GDB Online.}

\begin{frame}{Peta materi UTS}
\framesubtitle{Pertemuan 1 sampai 6}
\small
\begin{tabular}{@{}>{\raggedright\arraybackslash}p{124pt}>{\raggedright\arraybackslash}p{345pt}>{\raggedright\arraybackslash}p{393pt}@{}}
\kepala{Pertemuan} & \kepala{Konsep kunci} & \kepala{Jebakan yang sering muncul}\\
P1 & struktur program, \texttt{cout}, escape & garis miring tunggal, lupa titik koma\\
\barisgenap P2 & tipe data, \texttt{cin}, \texttt{getline} & pembagian bulat, \texttt{getline} sesudah \texttt{cin}\\
P3 & \texttt{\%}, precedence, \texttt{++}, logika & urutan operasi, \texttt{x++} vs \texttt{++x}\\
\barisgenap P4 & \texttt{if}, \texttt{else if}, \texttt{switch} & \texttt{=} vs \texttt{==}, urutan syarat, \texttt{break}\\
P5 & \texttt{for}, \texttt{while}, \texttt{do-while} & akumulator tidak diinisialisasi, off-by-one\\
\barisgenap P6 & bersarang, \texttt{break}, \texttt{continue} & batas dalam, rumus pola\\
\garisdasar
\end{tabular}
\end{frame}

\begin{frame}{Strategi soal A}
\framesubtitle{Menulis program}
\begin{columns}[T]
\begin{column}{0.31\textwidth}
\begin{Langkah}{1}{Rancang dulu}
Tulis masukan, proses, keluaran, dan data uji di kertas sebelum mengetik.
\end{Langkah}
\end{column}
\begin{column}{0.31\textwidth}
\begin{Langkah}{2}{Bangun bertahap}
Baca masukan, tampilkan, kompilasi. Tambah satu bagian, kompilasi lagi.
\end{Langkah}
\end{column}
\begin{column}{0.31\textwidth}
\begin{Langkah}{3}{Uji batas}
Jalankan dengan data uji biasa dan nilai batas. Cocokkan format keluaran.
\end{Langkah}
\end{column}
\end{columns}
\vspace{8pt}
Program yang sebagian benar dan terkompilasi mendapat nilai lebih baik daripada program lengkap yang tidak terkompilasi.
\end{frame}

\begin{frame}{Strategi soal B}
\framesubtitle{Tracing dan penjelasan}
\begin{columns}[T]
\begin{column}{0.31\textwidth}
\begin{Langkah}{1}{Buat tabel}
Satu kolom per variabel, satu baris per perubahan nilai.
\end{Langkah}
\end{column}
\begin{column}{0.31\textwidth}
\begin{Langkah}{2}{Ikuti persis}
Kerjakan sesuai aturan C++: pembagian bulat, precedence, \texttt{break}.
\end{Langkah}
\end{column}
\begin{column}{0.31\textwidth}
\begin{Langkah}{3}{Jelaskan}
Pakai istilah tepat: pembagian bulat, nilai batas, akumulator, cabang.
\end{Langkah}
\end{column}
\end{columns}
\vspace{8pt}
Jawaban B tanpa langkah hanya mendapat sebagian poin, walaupun hasil akhirnya benar.
\end{frame}

\begin{frame}[fragile]{Contoh tracing gabungan}
\framesubtitle{Soal B}
\begin{columns}[T]
\begin{column}{0.55\textwidth}
\begin{Kode}[]{gabungan.cpp}
int n = 6, hasil = 0;
for (int i = 1; i <= n; i++) {
    if (i % 2 == 0) {
        hasil += i;
    } else if (i % 3 == 0) {
        hasil -= i;
    }
}
cout << hasil << endl;
\end{Kode}
\end{column}
\begin{column}{0.41\textwidth}
\only<1>{\begin{TebakDulu}
{\ITERAKickerFont\fontsize{10pt}{12pt}\selectfont\color{ITERAAlert}TEBAK DULU}\par\vspace{4pt}
Sebelum dijalankan, tulis keluarannya di kertas.
\end{TebakDulu}
}\only<2>{\KeluaranFile[]{keluaran/P07/k001.txt}
\begin{Catatan}
i = 2, 4, 6 menambah. i = 3 mengurangi. i = 6 genap, jadi cabang kedua tidak diperiksa. Hasil: 2 + 4 + 6 − 3 = 9.
\end{Catatan}
}
\end{column}
\end{columns}
\note<1>{Minta mahasiswa menebak keluaran sebelum membuka slide berikutnya. Tanyakan satu atau dua tebakan yang berbeda, lalu buktikan.}
\end{frame}

\begin{frame}[fragile]{Memecah digit}
\framesubtitle{Pola yang sering keluar}
\begin{columns}[T]
\begin{column}{0.55\textwidth}
\begin{Kode}[]{digit.cpp}
int n = 4096, jumlah = 0;
while (n > 0) {
    jumlah += n % 10;
    n /= 10;
}
cout << jumlah << endl;
\end{Kode}
\end{column}
\begin{column}{0.41\textwidth}
\only<1>{\begin{TebakDulu}
{\ITERAKickerFont\fontsize{10pt}{12pt}\selectfont\color{ITERAAlert}TEBAK DULU}\par\vspace{4pt}
Sebelum dijalankan, tulis keluarannya di kertas.
\end{TebakDulu}
}\only<2>{\KeluaranFile[]{keluaran/P07/k002.txt}
\begin{Catatan}
\texttt{\% 10} mengambil digit terakhir, \texttt{/ 10} membuangnya. Pola ini muncul di banyak soal.
\end{Catatan}
}
\end{column}
\end{columns}
\note<1>{Minta mahasiswa menebak keluaran sebelum membuka slide berikutnya. Tanyakan satu atau dua tebakan yang berbeda, lalu buktikan.}
\end{frame}

\begin{frame}[fragile]{Tebak keluarannya}
\framesubtitle{Latihan menjelaskan, 3 menit}
\begin{columns}[T]
\begin{column}{0.56\textwidth}
\begin{Kode}[]{tebak.cpp}
int a = 5, b = 0;
while (a > 0) {
    if (a % 2 == 1) {
        b += a;
    } else {
        b -= 1;
    }
    a -= 2;
}
cout << a << " " << b << endl;
\end{Kode}
\end{column}
\begin{column}{0.40\textwidth}
\begin{Catatan}
Tulis keluarannya di kertas, karakter demi karakter. Jangan dijalankan dulu.
\end{Catatan}
\end{column}
\end{columns}
\note{Contoh soal Bagian B. Beri waktu 3 menit sebelum slide jawaban.}
\end{frame}

\begin{frame}[fragile]{Tebak keluarannya: jawaban}
\framesubtitle{Latihan menjelaskan}
\begin{columns}[T]
\begin{column}{0.56\textwidth}
\begin{Kode}[]{tebak.cpp}
int a = 5, b = 0;
while (a > 0) {
    if (a % 2 == 1) {
        b += a;
    } else {
        b -= 1;
    }
    a -= 2;
}
cout << a << " " << b << endl;
\end{Kode}
\end{column}
\begin{column}{0.40\textwidth}
\begin{Keluaran}[]
-1 9
\end{Keluaran}
\begin{Catatan}
\texttt{a} bernilai 5, 3, 1, semuanya ganjil, jadi \texttt{b} = 5 + 3 + 1 = 9. Sesudah itu \texttt{a} = −1 dan perulangan berhenti.
\end{Catatan}
\end{column}
\end{columns}
\note{Tanyakan jawaban yang paling sering salah, lalu telusuri bersama.}
\end{frame}

\begin{frame}{Error yang sering muncul minggu ini}
\framesubtitle{Pesan diambil langsung dari g++}
\footnotesize
\begin{tabular}{@{}>{\raggedright\arraybackslash}p{208pt}>{\raggedright\arraybackslash}p{256pt}>{\raggedright\arraybackslash}p{398pt}@{}}
\kepala{Kesalahan} & \kepala{Contoh} & \kepala{Pesan pertama dari g++}\\
Lupa titik koma & \texttt{cout << x} & \texttt{error: expected ';' before 'return'}\\
\barisgenap Variabel belum dideklarasikan & \texttt{total += x;} & \texttt{error: 'total' was not declared in this scope}\\
= di dalam kondisi & \texttt{if (a = b)} & \texttt{warning: suggest parentheses around assignment used as truth value}\\
\barisgenap break di luar perulangan & \texttt{break; di if} & \texttt{error: break statement not within loop or switch}\\
Kurung kurawal kurang & \texttt{while \{ if \{ \}} & \texttt{error: expected '\}' at end of input}\\
\garisdasar
\end{tabular}
\vspace{10pt}

{\small Pesan di tabel ini dihasilkan dengan g++ dan opsi -Wall, sama seperti di cek.sh.}
\note{Minta mahasiswa memotret slide ini.}
\end{frame}

\Bagian{4}{Hands-on bertingkat}{45 menit, dari dasar sampai tantangan}
\note{Mahasiswa membuka folder 02 Hands-on/P7 - Review UTS - Hands-on.}

\begin{frame}{Peta hands-on}
\framesubtitle{Folder 02 Hands-on/P7 - Review UTS - Hands-on}
\small
\begin{tabular}{@{}>{\raggedright\arraybackslash}p{36pt}>{\raggedright\arraybackslash}p{267pt}>{\raggedright\arraybackslash}p{110pt}>{\raggedright\arraybackslash}p{83pt}>{\raggedright\arraybackslash}p{341pt}@{}}
\kepala{No} & \kepala{File} & \kepala{Tingkat} & \kepala{Waktu} & \kepala{Yang dilatih}\\
1 & \texttt{handson1-parkir.cpp} & Dasar & 20' & masukan, \texttt{/} dan \texttt{\%}, \texttt{else if}\\
\barisgenap 2 & \texttt{handson2-digit.cpp} & Menengah & 15' & \texttt{while} dengan pola digit\\
3 & \texttt{handson3-telusur.cpp} & Tantangan & 10' & trace table bersarang dengan \texttt{continue} dan \texttt{break}\\
\barisgenap + & \texttt{bonus-menu.cpp} & Pengayaan & sisa & menu \texttt{do-while}, validasi\\
\garisdasar
\end{tabular}
\vspace{10pt}

Kerjakan berurutan. Cocokkan keluaran dengan folder \texttt{expected/}; latihan yang membaca masukan memakai data di folder \texttt{input/}.\note{Saat berkeliling, minta mahasiswa yang mengerjakan latihan tantangan menjelaskan blok PENJELASAN mereka.}
\end{frame}

\begin{frame}{Cara mengecek dan aturannya}
\framesubtitle{Hands-on}
\begin{columns}[T]
\begin{column}{0.47\textwidth}
\begin{Blok}{Di GDB Online}
Salin isi file latihan, lengkapi TODO, klik Run. Bila program membaca masukan, ketik isi file di \texttt{input/}. Bandingkan dengan \texttt{expected/}.
\end{Blok}
\end{column}
\begin{column}{0.47\textwidth}
\begin{BlokGelap}{Di komputer sendiri}
Jalankan \texttt{bash cek.sh}. Skrip mengompilasi dengan \texttt{-Wall -Werror}, memberi masukan dari \texttt{input/}, lalu mencocokkan keluaran.
\end{BlokGelap}
\end{column}
\end{columns}
\vspace{6pt}
\begin{Catatan}
AI boleh untuk memahami pesan error, tidak untuk meminta solusi utuh. Exit Ticket dikerjakan tanpa bantuan apa pun.
\end{Catatan}
\note{Hands-on tidak dinilai langsung; yang dinilai Exit Ticket dan tugas.}
\end{frame}

\Bagian{5}{Exit Ticket dan tugas}{25 menit terakhir, lalu pekerjaan rumah}
\note{Mulai Exit Ticket tepat waktu.}

\begin{frame}{Exit Ticket 7}
\framesubtitle{Individual, 25 menit, tanpa AI dan internet}
\small
\begin{tabular}{@{}>{\raggedright\arraybackslash}p{60pt}>{\raggedright\arraybackslash}p{560pt}>{\raggedright\arraybackslash}p{80pt}>{\raggedright\arraybackslash}p{150pt}@{}}
\kepala{Soal} & \kepala{Isi} & \kepala{Poin} & \kepala{CPMK}\\
A1 & Tulis program singkat yang memakai masukan, percabangan, dan perulangan & 25 & CPMK0607\\
\barisgenap A2 & Perbaiki program agar keluarannya sesuai contoh & 25 & CPMK0607\\
B1 & Buat trace table untuk perulangan bersarang dengan \texttt{continue} & 20 & CPMK0608\\
\barisgenap B2 & Jelaskan mengapa dua cabang tidak dijalankan bersamaan pada sebuah rantai \texttt{else if} & 20 & CPMK0608\\
B3 & Sebutkan materi pertemuan yang dipakai sebuah soal dan alasannya & 10 & CPMK0608\\
\garisdasar
\end{tabular}
\vspace{10pt}

{\small Soal A dinilai dari keluaran yang tepat sama. Soal B dinilai dari ketepatan dan alasan. Pelanggaran aturan membuat nilai Exit Ticket ini 0.}
\note{Soal lengkap dibagikan terpisah dan tidak disimpan di repo publik.}
\end{frame}

\begin{frame}{Sebelum Pertemuan 8}
\framesubtitle{Persiapan}
\begin{columns}[T]
\begin{column}{0.31\textwidth}
\begin{Langkah}{1}{Selesaikan}
Hands-on yang belum lulus \texttt{cek.sh}, terutama latihan tantangan.
\end{Langkah}
\end{column}
\begin{column}{0.31\textwidth}
\begin{Langkah}{2}{Latih}
Tidak ada tugas minggu ini. Kerjakan ulang ketiga simulasi dengan batas waktu ujian.
\end{Langkah}
\end{column}
\begin{column}{0.31\textwidth}
\begin{Langkah}{3}{Baca}
Modul Belajar 7 untuk mengulang, lalu bagian awal modul berikutnya.
\end{Langkah}
\end{column}
\end{columns}
\vspace{10pt}
Pertemuan 8 adalah UTS. Kisi-kisi, tata tertib, dan contoh soal ada di slide P8 dan Modul Belajar 8.\note{Tanyakan satu hal yang paling membingungkan hari ini untuk dibuka di review minggu depan.}
\end{frame}

\SlidePenutup{Terima kasih}{Sampai jumpa di UTS. Kerjakan contoh soal P8 dengan batas waktu.}
\note{Slide penutup.}
\end{document}
"
\documentclass[t]{beamer}
\usetheme{ITERA}
% Mode catatan pembicara: aktif bila \catatan didefinisikan sebelum \input.
\ifdefined\catatan
  \setbeameroption{show only notes}
  \setbeamertemplate{note page}{\vspace*{20pt}\hspace*{4pt}\begin{minipage}{900pt}
    {\ITERAKickerFont\fontsize{11pt}{14pt}\selectfont\color{ITERAGoldText}CATATAN PEMBICARA \ \ |\ \ SLIDE \insertframenumber}\par\vspace{4pt}
    {\ITERADisplay\bfseries\fontsize{22pt}{28pt}\selectfont\color{ITERAStructure}\insertframetitle}\par\vspace{12pt}
    \fontsize{15pt}{23pt}\selectfont\insertnote\end{minipage}}
\fi
\tikzset{fase/.style={draw=none,fill=#1,rounded corners=3pt,minimum width=158pt,minimum height=118pt,text width=146pt,align=center}}

\renewcommand\ITERAmeeting{Pertemuan 7}
\begin{document}
\SlideJudul{IF25-11002 PRAKTIKUM PEMROGRAMAN}{Pertemuan 7\\Review dan Simulasi UTS}{Merangkai materi Pertemuan 1 sampai 6 dan berlatih dengan soal bergaya UTS}{Tim Pengajar}{Tahun 2026. Sejajar dengan Pertemuan 7 Algoritma Pemrograman: Review dan Simulasi UTS.}
\note{Buka dengan menanyakan satu hal dari pertemuan lalu yang masih membingungkan.}

\begin{frame}{Dua mata kuliah, satu minggu}
\framesubtitle{Sebelum mulai}
Topik minggu ini sama di dua mata kuliah, dengan peran yang berbeda.
\vspace{10pt}

\begin{columns}[T]
\begin{column}{0.47\textwidth}
\begin{Blok}{Algoritma: Review dan Simulasi UTS}
Review pseudocode, flowchart, dan trace table untuk materi pertemuan 1 sampai 6.
\end{Blok}
\end{column}
\begin{column}{0.47\textwidth}
\begin{BlokGelap}{Praktikum: Review UTS}
Review C++ pertemuan 1 sampai 6 dan simulasi soal A (kode) dan soal B (tracing dan penjelasan).
\end{BlokGelap}
\end{column}
\end{columns}
\note{Tanyakan apa yang sudah dibahas di Algoritma minggu ini. Gunakan jawabannya sebagai jembatan ke praktikum.}
\end{frame}

\begin{frame}{Saat pulang nanti, kalian bisa}
\framesubtitle{Target hari ini}
\begin{columns}[T]
\begin{column}{0.47\textwidth}
\begin{Langkah}{1}{Menerapkan}
Menyelesaikan masalah yang menggabungkan masukan, ekspresi, percabangan, dan perulangan dalam satu program C++ dalam batas waktu ujian

{\footnotesize\color{ITERAInkMuted}Sub-CPMK 7.1, CPMK0607}
\end{Langkah}
\end{column}
\begin{column}{0.47\textwidth}
\begin{Langkah}{2}{Menjelaskan}
Menelusuri program yang menggabungkan percabangan dan perulangan bersarang, lalu menjelaskan setiap perubahan nilai dengan istilah yang tepat

{\footnotesize\color{ITERAInkMuted}Sub-CPMK 7.2, CPMK0608}
\end{Langkah}
\end{column}
\end{columns}
\vspace{8pt}
\begin{Catatan}
Dua kemampuan ini dinilai sama berat di Exit Ticket, tugas, UTS, dan UAS.
\end{Catatan}
\note{Bacakan target dengan pelan. Kata kuncinya menerapkan dan menjelaskan.}
\end{frame}

\begin{frame}{Ingat minggu lalu}
\framesubtitle{Review, 5 menit}
\begin{columns}[T]
\begin{column}{0.31\textwidth}
\begin{BlokGelap}{Pertanyaan 1}
Apa rumus banyak bintang pada piramida baris b?
\end{BlokGelap}
\end{column}
\begin{column}{0.31\textwidth}
\begin{BlokGelap}{Pertanyaan 2}
Apa beda \texttt{break} dan \texttt{continue} di perulangan bersarang?
\end{BlokGelap}
\end{column}
\begin{column}{0.31\textwidth}
\begin{BlokGelap}{Pertanyaan 3}
Bagaimana menghitung jam dari menit dengan pembulatan ke atas?
\end{BlokGelap}
\end{column}
\end{columns}
\vspace{8pt}
Jawab di kertas dulu, lalu cocokkan dengan teman sebelah.\note{Gunakan tiga pertanyaan ini sebagai diagnosis cepat. Catat topik yang paling banyak salah untuk dibahas di bagian materi.}
\end{frame}

\Bagian{1}{Masalah pembuka}{Satu soal, semua materi}
\note{Masalah pembuka 10 menit.}

\begin{frame}[fragile]{Satu soal, semua materi}
\framesubtitle{Masalah minggu ini}
\begin{columns}[T]
\begin{column}{0.55\textwidth}
{\small Soal UTS jarang menguji satu topik saja. Contohnya tarif parkir: membaca masukan, membulatkan menit ke jam dengan \texttt{/} dan \texttt{\%}, memilih tarif dengan \texttt{if}, dan mengulang untuk beberapa kendaraan.

\vspace{6pt}
Hari ini kita berlatih memecah soal seperti itu menjadi bagian-bagian yang sudah kalian kuasai.\par}
\vspace{6pt}
\begin{Catatan}
\textbf{Diskusi 2 menit.} Materi pertemuan mana saja yang dipakai soal ini? Bagian mana yang menurut kalian paling rawan salah?
\end{Catatan}
\end{column}
\begin{column}{0.41\textwidth}
\begin{Keluaran}[]
M 61   -> 2 jam, tarif 3000
B 180  -> 3 jam, tarif 11000
X 10   -> Jenis tidak dikenal
\end{Keluaran}
\end{column}
\end{columns}
\note{Tampung beberapa jawaban tanpa menilai benar salah.}
\end{frame}

\begin{frame}{Dari langkah ke instruksi}
\framesubtitle{Sambungan dengan Algoritma}
\begin{columns}[T]
\begin{column}{0.31\textwidth}
\begin{Langkah}{1}{Baca soal}
Garis bawahi masukan, keluaran, dan aturan.
\end{Langkah}
\end{column}
\begin{column}{0.31\textwidth}
\begin{Langkah}{2}{Pecah}
Petakan setiap aturan ke konstruksi C++ yang sudah dipelajari.
\end{Langkah}
\end{column}
\begin{column}{0.31\textwidth}
\begin{Langkah}{3}{Uji}
Siapkan data uji, termasuk nilai batas, sebelum menulis kode.
\end{Langkah}
\end{column}
\end{columns}
\vspace{10pt}
Langkah ini sama dengan yang dilatih di Algoritma minggu ini, hanya hasil akhirnya kode C++.\note{Hubungkan eksplisit ke Algoritma minggu ini.}
\end{frame}

\Bagian{2}{Coba dulu, baru dijelaskan}{25 menit eksperimen di GDB Online}
\note{Struggle 25 menit. Berkeliling, jangan menjelaskan materi dulu.}

\begin{frame}{Misi kalian}
\framesubtitle{Struggle, 25 menit}
\begin{columns}[T]
\begin{column}{0.31\textwidth}
\begin{Langkah}{1}{Pilih}
Kerjakan \texttt{handson1-parkir.cpp} tanpa bantuan selama 20 menit.
\end{Langkah}
\end{column}
\begin{column}{0.31\textwidth}
\begin{Langkah}{2}{Uji}
Uji dengan 60 dan 61 menit. Apakah hasilnya berbeda satu jam?
\end{Langkah}
\end{column}
\begin{column}{0.31\textwidth}
\begin{Langkah}{3}{Catat}
Tulis satu hal yang membuat kalian tersendat.
\end{Langkah}
\end{column}
\end{columns}
\vspace{8pt}
\begin{columns}[T]
\begin{column}{0.47\textwidth}
\begin{Blok}{Boleh}
Membuka catatan dan modul sendiri.
\end{Blok}
\end{column}
\begin{column}{0.47\textwidth}
\begin{BlokAlert}{Belum dulu}
Bertanya ke teman atau AI. Ini simulasi kondisi ujian.
\end{BlokAlert}
\end{column}
\end{columns}
\note{Catat kesalahan yang sering muncul untuk dibahas di debrief.}
\end{frame}

\begin{frame}{Apa yang kalian temukan?}
\framesubtitle{Debrief struggle}
\begin{columns}[T]
\begin{column}{0.31\textwidth}
\begin{BlokGelap}{Pertanyaan 1}
Bagian mana yang paling lama dikerjakan?
\end{BlokGelap}
\end{column}
\begin{column}{0.31\textwidth}
\begin{BlokGelap}{Pertanyaan 2}
Bagaimana kalian membulatkan menit ke atas?
\end{BlokGelap}
\end{column}
\begin{column}{0.31\textwidth}
\begin{BlokGelap}{Pertanyaan 3}
Data uji apa yang kalian pakai?
\end{BlokGelap}
\end{column}
\end{columns}
\vspace{10pt}
Jawaban kalian akan kita uji di bagian berikutnya.\note{Tulis jawaban di papan; buktikan atau patahkan di bagian materi.}
\end{frame}

\Bagian{3}{Merangkai pertemuan 1 sampai 6}{Peta konsep, strategi soal A, dan strategi soal B}
\note{Materi 40 menit. Tunjukkan setiap contoh langsung di GDB Online.}

\begin{frame}{Peta materi UTS}
\framesubtitle{Pertemuan 1 sampai 6}
\small
\begin{tabular}{@{}>{\raggedright\arraybackslash}p{124pt}>{\raggedright\arraybackslash}p{345pt}>{\raggedright\arraybackslash}p{393pt}@{}}
\kepala{Pertemuan} & \kepala{Konsep kunci} & \kepala{Jebakan yang sering muncul}\\
P1 & struktur program, \texttt{cout}, escape & garis miring tunggal, lupa titik koma\\
\barisgenap P2 & tipe data, \texttt{cin}, \texttt{getline} & pembagian bulat, \texttt{getline} sesudah \texttt{cin}\\
P3 & \texttt{\%}, precedence, \texttt{++}, logika & urutan operasi, \texttt{x++} vs \texttt{++x}\\
\barisgenap P4 & \texttt{if}, \texttt{else if}, \texttt{switch} & \texttt{=} vs \texttt{==}, urutan syarat, \texttt{break}\\
P5 & \texttt{for}, \texttt{while}, \texttt{do-while} & akumulator tidak diinisialisasi, off-by-one\\
\barisgenap P6 & bersarang, \texttt{break}, \texttt{continue} & batas dalam, rumus pola\\
\garisdasar
\end{tabular}
\end{frame}

\begin{frame}{Strategi soal A}
\framesubtitle{Menulis program}
\begin{columns}[T]
\begin{column}{0.31\textwidth}
\begin{Langkah}{1}{Rancang dulu}
Tulis masukan, proses, keluaran, dan data uji di kertas sebelum mengetik.
\end{Langkah}
\end{column}
\begin{column}{0.31\textwidth}
\begin{Langkah}{2}{Bangun bertahap}
Baca masukan, tampilkan, kompilasi. Tambah satu bagian, kompilasi lagi.
\end{Langkah}
\end{column}
\begin{column}{0.31\textwidth}
\begin{Langkah}{3}{Uji batas}
Jalankan dengan data uji biasa dan nilai batas. Cocokkan format keluaran.
\end{Langkah}
\end{column}
\end{columns}
\vspace{8pt}
Program yang sebagian benar dan terkompilasi mendapat nilai lebih baik daripada program lengkap yang tidak terkompilasi.
\end{frame}

\begin{frame}{Strategi soal B}
\framesubtitle{Tracing dan penjelasan}
\begin{columns}[T]
\begin{column}{0.31\textwidth}
\begin{Langkah}{1}{Buat tabel}
Satu kolom per variabel, satu baris per perubahan nilai.
\end{Langkah}
\end{column}
\begin{column}{0.31\textwidth}
\begin{Langkah}{2}{Ikuti persis}
Kerjakan sesuai aturan C++: pembagian bulat, precedence, \texttt{break}.
\end{Langkah}
\end{column}
\begin{column}{0.31\textwidth}
\begin{Langkah}{3}{Jelaskan}
Pakai istilah tepat: pembagian bulat, nilai batas, akumulator, cabang.
\end{Langkah}
\end{column}
\end{columns}
\vspace{8pt}
Jawaban B tanpa langkah hanya mendapat sebagian poin, walaupun hasil akhirnya benar.
\end{frame}

\begin{frame}[fragile]{Contoh tracing gabungan}
\framesubtitle{Soal B}
\begin{columns}[T]
\begin{column}{0.55\textwidth}
\begin{Kode}[]{gabungan.cpp}
int n = 6, hasil = 0;
for (int i = 1; i <= n; i++) {
    if (i % 2 == 0) {
        hasil += i;
    } else if (i % 3 == 0) {
        hasil -= i;
    }
}
cout << hasil << endl;
\end{Kode}
\end{column}
\begin{column}{0.41\textwidth}
\only<1>{\begin{TebakDulu}
{\ITERAKickerFont\fontsize{10pt}{12pt}\selectfont\color{ITERAAlert}TEBAK DULU}\par\vspace{4pt}
Sebelum dijalankan, tulis keluarannya di kertas.
\end{TebakDulu}
}\only<2>{\KeluaranFile[]{keluaran/P07/k001.txt}
\begin{Catatan}
i = 2, 4, 6 menambah. i = 3 mengurangi. i = 6 genap, jadi cabang kedua tidak diperiksa. Hasil: 2 + 4 + 6 − 3 = 9.
\end{Catatan}
}
\end{column}
\end{columns}
\note<1>{Minta mahasiswa menebak keluaran sebelum membuka slide berikutnya. Tanyakan satu atau dua tebakan yang berbeda, lalu buktikan.}
\end{frame}

\begin{frame}[fragile]{Memecah digit}
\framesubtitle{Pola yang sering keluar}
\begin{columns}[T]
\begin{column}{0.55\textwidth}
\begin{Kode}[]{digit.cpp}
int n = 4096, jumlah = 0;
while (n > 0) {
    jumlah += n % 10;
    n /= 10;
}
cout << jumlah << endl;
\end{Kode}
\end{column}
\begin{column}{0.41\textwidth}
\only<1>{\begin{TebakDulu}
{\ITERAKickerFont\fontsize{10pt}{12pt}\selectfont\color{ITERAAlert}TEBAK DULU}\par\vspace{4pt}
Sebelum dijalankan, tulis keluarannya di kertas.
\end{TebakDulu}
}\only<2>{\KeluaranFile[]{keluaran/P07/k002.txt}
\begin{Catatan}
\texttt{\% 10} mengambil digit terakhir, \texttt{/ 10} membuangnya. Pola ini muncul di banyak soal.
\end{Catatan}
}
\end{column}
\end{columns}
\note<1>{Minta mahasiswa menebak keluaran sebelum membuka slide berikutnya. Tanyakan satu atau dua tebakan yang berbeda, lalu buktikan.}
\end{frame}

\begin{frame}[fragile]{Tebak keluarannya}
\framesubtitle{Latihan menjelaskan, 3 menit}
\begin{columns}[T]
\begin{column}{0.56\textwidth}
\begin{Kode}[]{tebak.cpp}
int a = 5, b = 0;
while (a > 0) {
    if (a % 2 == 1) {
        b += a;
    } else {
        b -= 1;
    }
    a -= 2;
}
cout << a << " " << b << endl;
\end{Kode}
\end{column}
\begin{column}{0.40\textwidth}
\begin{Catatan}
Tulis keluarannya di kertas, karakter demi karakter. Jangan dijalankan dulu.
\end{Catatan}
\end{column}
\end{columns}
\note{Contoh soal Bagian B. Beri waktu 3 menit sebelum slide jawaban.}
\end{frame}

\begin{frame}[fragile]{Tebak keluarannya: jawaban}
\framesubtitle{Latihan menjelaskan}
\begin{columns}[T]
\begin{column}{0.56\textwidth}
\begin{Kode}[]{tebak.cpp}
int a = 5, b = 0;
while (a > 0) {
    if (a % 2 == 1) {
        b += a;
    } else {
        b -= 1;
    }
    a -= 2;
}
cout << a << " " << b << endl;
\end{Kode}
\end{column}
\begin{column}{0.40\textwidth}
\begin{Keluaran}[]
-1 9
\end{Keluaran}
\begin{Catatan}
\texttt{a} bernilai 5, 3, 1, semuanya ganjil, jadi \texttt{b} = 5 + 3 + 1 = 9. Sesudah itu \texttt{a} = −1 dan perulangan berhenti.
\end{Catatan}
\end{column}
\end{columns}
\note{Tanyakan jawaban yang paling sering salah, lalu telusuri bersama.}
\end{frame}

\begin{frame}{Error yang sering muncul minggu ini}
\framesubtitle{Pesan diambil langsung dari g++}
\footnotesize
\begin{tabular}{@{}>{\raggedright\arraybackslash}p{208pt}>{\raggedright\arraybackslash}p{256pt}>{\raggedright\arraybackslash}p{398pt}@{}}
\kepala{Kesalahan} & \kepala{Contoh} & \kepala{Pesan pertama dari g++}\\
Lupa titik koma & \texttt{cout << x} & \texttt{error: expected ';' before 'return'}\\
\barisgenap Variabel belum dideklarasikan & \texttt{total += x;} & \texttt{error: 'total' was not declared in this scope}\\
= di dalam kondisi & \texttt{if (a = b)} & \texttt{warning: suggest parentheses around assignment used as truth value}\\
\barisgenap break di luar perulangan & \texttt{break; di if} & \texttt{error: break statement not within loop or switch}\\
Kurung kurawal kurang & \texttt{while \{ if \{ \}} & \texttt{error: expected '\}' at end of input}\\
\garisdasar
\end{tabular}
\vspace{10pt}

{\small Pesan di tabel ini dihasilkan dengan g++ dan opsi -Wall, sama seperti di cek.sh.}
\note{Minta mahasiswa memotret slide ini.}
\end{frame}

\Bagian{4}{Hands-on bertingkat}{45 menit, dari dasar sampai tantangan}
\note{Mahasiswa membuka folder 02 Hands-on/P7 - Review UTS - Hands-on.}

\begin{frame}{Peta hands-on}
\framesubtitle{Folder 02 Hands-on/P7 - Review UTS - Hands-on}
\small
\begin{tabular}{@{}>{\raggedright\arraybackslash}p{36pt}>{\raggedright\arraybackslash}p{267pt}>{\raggedright\arraybackslash}p{110pt}>{\raggedright\arraybackslash}p{83pt}>{\raggedright\arraybackslash}p{341pt}@{}}
\kepala{No} & \kepala{File} & \kepala{Tingkat} & \kepala{Waktu} & \kepala{Yang dilatih}\\
1 & \texttt{handson1-parkir.cpp} & Dasar & 20' & masukan, \texttt{/} dan \texttt{\%}, \texttt{else if}\\
\barisgenap 2 & \texttt{handson2-digit.cpp} & Menengah & 15' & \texttt{while} dengan pola digit\\
3 & \texttt{handson3-telusur.cpp} & Tantangan & 10' & trace table bersarang dengan \texttt{continue} dan \texttt{break}\\
\barisgenap + & \texttt{bonus-menu.cpp} & Pengayaan & sisa & menu \texttt{do-while}, validasi\\
\garisdasar
\end{tabular}
\vspace{10pt}

Kerjakan berurutan. Cocokkan keluaran dengan folder \texttt{expected/}; latihan yang membaca masukan memakai data di folder \texttt{input/}.\note{Saat berkeliling, minta mahasiswa yang mengerjakan latihan tantangan menjelaskan blok PENJELASAN mereka.}
\end{frame}

\begin{frame}{Cara mengecek dan aturannya}
\framesubtitle{Hands-on}
\begin{columns}[T]
\begin{column}{0.47\textwidth}
\begin{Blok}{Di GDB Online}
Salin isi file latihan, lengkapi TODO, klik Run. Bila program membaca masukan, ketik isi file di \texttt{input/}. Bandingkan dengan \texttt{expected/}.
\end{Blok}
\end{column}
\begin{column}{0.47\textwidth}
\begin{BlokGelap}{Di komputer sendiri}
Jalankan \texttt{bash cek.sh}. Skrip mengompilasi dengan \texttt{-Wall -Werror}, memberi masukan dari \texttt{input/}, lalu mencocokkan keluaran.
\end{BlokGelap}
\end{column}
\end{columns}
\vspace{6pt}
\begin{Catatan}
AI boleh untuk memahami pesan error, tidak untuk meminta solusi utuh. Exit Ticket dikerjakan tanpa bantuan apa pun.
\end{Catatan}
\note{Hands-on tidak dinilai langsung; yang dinilai Exit Ticket dan tugas.}
\end{frame}

\Bagian{5}{Exit Ticket dan tugas}{25 menit terakhir, lalu pekerjaan rumah}
\note{Mulai Exit Ticket tepat waktu.}

\begin{frame}{Exit Ticket 7}
\framesubtitle{Individual, 25 menit, tanpa AI dan internet}
\small
\begin{tabular}{@{}>{\raggedright\arraybackslash}p{60pt}>{\raggedright\arraybackslash}p{560pt}>{\raggedright\arraybackslash}p{80pt}>{\raggedright\arraybackslash}p{150pt}@{}}
\kepala{Soal} & \kepala{Isi} & \kepala{Poin} & \kepala{CPMK}\\
A1 & Tulis program singkat yang memakai masukan, percabangan, dan perulangan & 25 & CPMK0607\\
\barisgenap A2 & Perbaiki program agar keluarannya sesuai contoh & 25 & CPMK0607\\
B1 & Buat trace table untuk perulangan bersarang dengan \texttt{continue} & 20 & CPMK0608\\
\barisgenap B2 & Jelaskan mengapa dua cabang tidak dijalankan bersamaan pada sebuah rantai \texttt{else if} & 20 & CPMK0608\\
B3 & Sebutkan materi pertemuan yang dipakai sebuah soal dan alasannya & 10 & CPMK0608\\
\garisdasar
\end{tabular}
\vspace{10pt}

{\small Soal A dinilai dari keluaran yang tepat sama. Soal B dinilai dari ketepatan dan alasan. Pelanggaran aturan membuat nilai Exit Ticket ini 0.}
\note{Soal lengkap dibagikan terpisah dan tidak disimpan di repo publik.}
\end{frame}

\begin{frame}{Sebelum Pertemuan 8}
\framesubtitle{Persiapan}
\begin{columns}[T]
\begin{column}{0.31\textwidth}
\begin{Langkah}{1}{Selesaikan}
Hands-on yang belum lulus \texttt{cek.sh}, terutama latihan tantangan.
\end{Langkah}
\end{column}
\begin{column}{0.31\textwidth}
\begin{Langkah}{2}{Latih}
Tidak ada tugas minggu ini. Kerjakan ulang ketiga simulasi dengan batas waktu ujian.
\end{Langkah}
\end{column}
\begin{column}{0.31\textwidth}
\begin{Langkah}{3}{Baca}
Modul Belajar 7 untuk mengulang, lalu bagian awal modul berikutnya.
\end{Langkah}
\end{column}
\end{columns}
\vspace{10pt}
Pertemuan 8 adalah UTS. Kisi-kisi, tata tertib, dan contoh soal ada di slide P8 dan Modul Belajar 8.\note{Tanyakan satu hal yang paling membingungkan hari ini untuk dibuka di review minggu depan.}
\end{frame}

\SlidePenutup{Terima kasih}{Sampai jumpa di UTS. Kerjakan contoh soal P8 dengan batas waktu.}
\note{Slide penutup.}
\end{document}
